\documentclass[instructions.tex]{subfiles}
\begin{document}
 
\chapter{Des danses et des bals.}
 
\primo Dans les villes, les bals et les danses se font avec plus d'éclat, et ordinairement avec plus de scandale; dans les campagnes, avec moins d'appareil, mais toujours avec danger.

Les bals qui se font en masque sont les plus condamnables: S'il n'est pas permis d'être déguisé dans ses paroles, sera-t-il permis de déguiser sous le masque sa profession, son nom, son sexe, sa personne, et, sous une figure ridicule, se faire un divertissement de ne paroître ni chrétien ni homme?

Saint Cyprien, saint Augustin, et saint Thomas enseignent que le bal fait injure à Dieu, parce qu'il réforme ce que Dieu a formé. Le masque qui défigure l'ouvrage de Dieu, est-il moins injurieux au Créateur? Néanmoins la folie du monde va jusqu'à dire que tout homme en masque est respectable. C'est donc à dire qu'un fat, un impudent, un homme infâme, a droit sous le masque, d'aller de pair avec un magistrat, avec un prince. Disons plutôt qu'il n'est rien de plus méprisable qu'un masque, et que les sages du paganisme ont regardé avec raison ces extravagantes mascarades comme indignes de l'homme, et que la religion les condamne.

\secundo Quant aux autres danses, on doit dire, avec saint François de Sales, que si la danse est une action de soi indifférente, les circonstances de la danse, surtout avec des personnes de différent sexe, la rendent toujours dangereuse, et rarement innocente. L'âme y reçoit des atteintes par les pensées, par les regards, par l'air de dissolution et de liberté qui règne dans ces sortes d'assemblées.

Quoiqu'en pense le monde, il ne doit pas ici nous servir de règle. Il ne faut qu'une impiété de dire que Dieu approuve le bal et la danse, lui qui nous ordonne de réprimer jusqu'à nos regards, qui nous défend même d'aimer la compagnie d'une danseuse. N'arrêtez pas vos regards sur une femme volage, dit le Seigneur dans les Livres saints. Gardez-vous de fréquenter une danseuse, ne l'écoutez même pas, crainte de périr par ses attraits. N'en voyez pas la vue sur une vierge... sur une femme parée... Leur beauté est un écueil qui en a fait périr plusieurs\footnote{Eccles. 9.}. C'est Dieu qui parle ainsi, et qui nous avertit du danger.

Nous pourrions citer l'autorité des SS. Pères, qui, tous, sans exception, réprouvent les danses profanes; mais nous nous contentons de rapporter le témoignage d'un seigneur de la cour, et grand capitaine. Voici comment il en parle dans un discours imprimé, qu'il adresse à ses enfans:

J'ai toujours craint les bals dangereux; ce n'a pas été seulement ma raison qui me l'a fait croire, ç'a encore été mon expérience. Quoique le témoignage des saints Pères, qui les condamnent soit bien fort, je tiens que, sur ce chapitre, le témoignage d'un courtisan doit être de plus grand poids. Je sais bien qu'il y a des gens qui y courent moins de hasards que d'autres, cependant les tempéramens les plus froids s'y échauffent. Ce ne sont d'ordinaire que des jeunes gens qui composent ces sortes d'assemblées, lesquels ont assez de peine à résister aux tentations dans la solitude; à plus forte raison dans ces lieux-là, où la beauté des objets, les illuminations, les violons, l'agitation de la danse, réveilleront les passions dans une chair très... Ainsi, je tiens qu'il ne faut point aller au bal; quand on est invité, et je crois que les directeurs forment leur devoir, s'ils exigent de ceux dont ils gouvernent la conscience, qu'ils n'y aillent jamais. Le comte de Bussy.

Ce Seigneur avoit éprouvé que le bal et la danse sont dangereux à l'innocence. Combien de gens ont éprouvé qu'ils leur ont été funestes!

\section{Résolutions}

\primo Je m'interdis tout bal public.

\secundo Si les bienséances m'obligent quelquefois de me trouver dans des bals particuliers, où peu de personnes et seulement des parens ou d'autres gens de choix seroient invités, je ne m'y rendrai qu'après avoir médité sérieusement quelque grande vérité de la religion, avoir beaucoup prié, m'être pénétré vivement de la présence de Dieu, du néant des plaisirs du monde, des dangers qu'on y court de perdre son âme; encore ne m'y déciderai-je que le plus rarement que je pourrai, et qu'avec la permission expresse d'un confesseur sage et expérimenté. 

\prayer{ô mon Dieu! de tous les plaisirs qui conviennent peu à la mortification chrétienne, et qui nuisent presque toujours au salut.}
 
\end{document}
